\section{A polynomial bound for a single cylinder arc}
\label{sec:one-arc}

The marked-pair estimate uses a sharp fluctuation determinant.  The
single-arc application needs only a polynomial upper bound, and this
weaker objective permits a larger parameter range.  In particular,
the shorter wall may grow arbitrarily slowly compared with the longer
one.  We prove the contour identity first, including its finite spin
and residue calculation, and then bound its absolute integral by the
real reference comparison.  No imaginary displacement of particles is
used in that comparison.

Retain the cyclic order $[p,X,p',Y]$, the magnetic factors $P_j$ and
$\mathcal T$, and the constants $\beta=\pi/4$ and $\gamma=4/3$.
There are $s=|X|$ and $b=|Y|$ ordinary sites, and $N=s+b+2$ is even.
Both marks and all ordinary sites have spectral heights in
$\{0,\beta\}$; every $Y$ site has height zero.  The variable
$\rho_{\rm v}=1/\sqrt{2+\sqrt2}$ denotes the vertex activity,
whereas $\rho(u)$ continues to denote the centering probability
density.  All arc masses are unsigned critical masses with the
previously specified port convention.

\subsection{The one-arc contour identity}

Let \(A(p,p')\) be the mass of self-avoiding cylinder arcs departing
on the west, or ket, side of a marked port \(p\), and arriving on the
east, or bra, side of a marked port \(p'\).  The arc has no other visit
to either port, and their unused half-edges are vacant.  Prescribe spin
\(+\) on the ket at \(p\), spin \(+\) on the bra at \(p'\), and
vacancy on the opposite sides, with ordinary-site twist
\(k^{-2\sum_{i\in X}h_i}\).  In the bilinear spin notation this is
\[
 \mathcal A(p,p')=
 \sum_{\mathbf h\in\{0,\pm1\}^{X\cup Y}}
 k^{-2\sum_{i\in X}h_i}
 U_{+,\mathbf h_X,0,\mathbf h_Y}
 U^*_{0,\mathbf h_X,+,\mathbf h_Y}.
\]

\begin{lemma}[One-arc contraction]\label{09_fast_paths:one-arc-contraction}
There is a constant $\omega$ with $|\omega|=1$, independent of the
arc, such that $\mathcal A(p,p')=\omega A(p,p')$.
\end{lemma}
\begin{proof}
A closed loop disjoint from the marked arc cannot separate its endpoints,
so its two orientations cancel as in the marked-pair calculation.  Order
the vertical matching from \(p\).  The directed arc has westward initial
and final tangents.  With argument cut below \(p\), the angle change of
\(\Gamma-p\) is \(-\pi/2\); with argument cut upward, that of
\(p'-\Gamma\) is \(-3\pi/2-2\pi F_X\).  The chord-angle identity
therefore cancels the twist against the cup phases, leaving a constant
phase.  The summability and projection of the half-cylinder states are
the same as for the marked-pair contraction.
\end{proof}

There are \(n=s+b\) ordinary sites and two marks, with \(N=n+2\) even.
In the finite calculation below \(t=+1,-1\) indexes the groups
\(X,Y\); it is not the fraction of upper ordinary sites. Write
\(p_t\) for the logarithmic coordinate of the mark preceding group
\(t\), rather than for an activity.  Define two probe factors
\begin{align*}
 K_t&=\frac{P_{-3t}(y_t-p_{-t})}
 {P_{2t}(y_t-p_t)P_0(y_t-p_t)},\\
 Q_t&=\prod_{a\text{ ordinary}}P_{s_a-t}(a-y_t)^{-1}
 \prod_{v\text{ elementary}}
 \frac{P_{s_v-t+1}(v-y_t)P_{s_v-t-1}(v-y_t)}
 {P_{s_v-t}(v-y_t)}.
\end{align*}
The signs \(s_a,s_v\) specify groups; a double contributes both its
elementary particles.  Integrate \(y_+\) around the strip with height
interval \((-0.8,3.34)\), and \(y_-\) around
\((-2.72,1.23)\), in units of \(\beta\), with measure
\(dy_t/(2\pi i)\) and positive orientation: lower horizontal integral
minus upper horizontal integral.  Ordinary sites have heights 0 or 1,
and all ordinary minus sites have height 0.  Initially the main lines
have heights \(1/2,0\), the far singleton lines have heights
\(25/8,-2.35\), and doubles have these far endpoints and their endpoints
one step inward.

\begin{lemma}[One-arc gas identity]\label{09_fast_paths:one-arc-identity}
Use the ordinary magnetic activity coefficients, including signs, but
elementary particle totals
\[
                  (s+m-1,b-m),\qquad m\in\mathbb Z,
\]
restricted to nonnegative values.  The preceding contraction multiplied
by \(\widetilde h_N^2\) is the sum of these gas integrals with factor
\(\mathcal T K_+K_-Q_+Q_-\), multiplied by \((2d)^n\) and sector
scalars of uniformly bounded modulus.  A scalar may depend on
\(m,n,s,b\), but not on the chosen pole sites or strings in its sector.
\end{lemma}
\begin{proof}
We first compute residues and then compare every exceptional spin factor.
Keep real site offsets generic and distinct until taking confluence.
Moving main contours to the far lines produces ordinary strings of
length 0, 1 or 2; the cancellation excluding longer strings is unchanged.
All singleton poles of \(Q_t\) lie outside the swept regions.  For a
double, multiplication of its elementary factors cancels the relevant
particle-dependent denominator before moving contours.  The additional
real growth is bounded by
\[
 \sum_{t=\pm1}\left\{
 \frac12\sum_v|\Re v-\Re y_t|
               -\frac{n+1}{2}|\Re y_t|\right\}.
\]
There are \(n-1\) ordinary elementary particles.  Up to a constant
depending on the counts and fixed site offsets, their ordinary gas
exponent is at most
\[
 -n\sum_v|\Re v|+\sum_{\{v,w\}}|\Re v-\Re w|
 \le -2\sum_v|\Re v|,
\]
where the second sum is over unordered elementary pairs.  Omitting
an internal double pair does not change this bound, since its real
difference is zero.  By the same triangle inequality, the displayed
probe contribution is at most
\[
 \sum_v|\Re v|-\sum_{t=\pm1}|\Re y_t|.
\]
Adding the two bounds gives simultaneous exponential decay in all
particles and both probes, justifying closure of the strips.

For a group-\(z\) site with string length \(a\), its site and string
contribution to \(Q_t\), at argument site minus probe, is
\[
\begin{array}{c|ccc}
 a&0&1&2\\ \hline
 &1/P_J&P_{J+2z}/P_{J+z}&P_{J+3z}
\end{array},\qquad J=z-t.
\]
The probe residues are therefore type \(\mathrm A\), at \(y_t=p_t\),
or type \(\mathrm B_a\), at \(y_t=a_j+ita\beta\), where \(a_j\)
is an ordinary site of group \(t\) with string length \(a=0,1\).
There are no other poles in the indicated strips.

Expand in the full-group biorthogonal bases \(D_l,D_r^*\), including
each initial mark, using Lemma~\ref{lem:06-dual-covectors}.  The summand
contains \(D_lU,D_r^*U^*\) divided by their two norms; the remaining
\(D_l^*,D_r\) meet the prescribed operator.  An ordinary permutation
in either group leaves the summand invariant: braid transport multiplies
vectors and their norms by the same scalars, and the prescribed operator
intertwines the ordinary braids.  Place an ordinary site last.  Its
untouched last roots exclude \(r_j<l_j\).  Projection gives
\(\sum l=m\) in both groups and \(\sum(l-r)=1\) in each group.
The plus mark has \(l_p=1\), while the minus mark has \(r_p=-1\),
by their leading occupied roots.  Thus the possibilities are precisely
\[
\begin{array}{c|cc|c}
 &t=+1&t=-1&\text{ordinary exceptions}\\ \hline
\mathrm A&(l_p,r_p)=(1,0)&(l_p,r_p)=(0,-1)&\text{none}\\
\mathrm B_a&(l_p,r_p)=(1,-1)&(l_p,r_p)=(1,-1)&
 r_j=l_j+1\text{ at one site }j.
\end{array}
\]
Other labels match.  The string length at any ordinary site is
\(\min(1+t l_j,1+t r_j)\), giving exactly the claimed particle totals.

Sort nonexceptional spectators first.  On the plus group's \(D_l^*\)
the occupied mark traverses minus roots with saturated spins: each
emerged minus forces nondecrease, and the initial spin is maximal.
These crossings supply their \(F\)-factors.  In type \(\mathrm A\)
the vacancy on \(D_r\) splits to \(+,-\), since its leading root is
plus; the second strand is again saturated.  In type \(\mathrm B\)
it starts vacant and ends minus while crossing plus roots, so exactly
one crossing contributes the unit producing an emerged plus.  This
must be the leading plus root of site \(j\).  If \(l_j=-1\), the
actual spin on \(D_l^*\) is minus, forcing the real split to vacancy,
minus: a minus could not emerge plus at the leading root.  If
\(l_j=0\), the actual spin must be vacant, and the starred split there
is minus, plus.  Their first relevant roots are untouched when the
ordinary sites are sorted before inserting the mark.  All remaining
exceptional strands then cross spectators with saturated spins.
Consequently there is exactly one output cup \(U\) in place of one
\(F\)-factor.  For the minus group exchange starred and real lists
and negate relative spins; the physical cup phases agree.  Spectator
overlap gives the matching norm.

Here are explicit rules for the product comparison.  Recall
\[
 \mathcal R(\pm1)=\{(\pm1,0)\},\qquad
 \mathcal R(0)=\{(1,-1),(-1,1)\}.
\]
For ordered sites, evaluate factors at first coordinate minus second.
Roots \((e,g),(f,h)\) shift that difference by \(g-h\).  Put \(u=1\)
for the \(l\)-half, \(u=-1\) for the \(r\)-half.  Equal groups and
equal signs contribute \(P_{2u}P_{3u}\); opposite groups and opposite
signs contribute \(P_{-2e}P_{-3e}\).  For equal groups and opposite
signs an uncancelled norm contributes
\[
                       \frac{P_{-2e}P_{-3e}}{P_0P_e}.
\]
Include this once, from the \(l\)-list, for matching ordinary
spectators.  If the earlier site is exceptional, include it in each
half only for \(e=-u\), as the saturated routes require.  In each
half divide by \(2\cosh((\cdot)/2)\) at shift zero if exactly one
site label is zero, or at shifts \(-1,0,1\) if both are zero.  Finally
the special cup contributes \((-2itd)/P_{-ta}\) at mark minus changed
site.  These rules account for both occupied cyclic blocks and both
centered merge quotients.

The ordinary one-site baseline is \((2c)^{(a-1)^2}\); use the same
baseline at the changed site.  The additional exceptional factors are
\[
 C_0=\frac{\rho_{\rm v}(2c)^{3/2}}{2dF(2\lambda)}\quad(\mathrm A),
 \qquad C_0\frac{2c}{(2c)^{(a-1)^2}}\quad(\mathrm B_a).
\]
Each split supplies centered divisor \(2d\), one internal split norm
is uncancelled, and the forced split coefficient is \(\rho_{\rm v}\), up to
phase.  Let \(z_{l,t},z_{r,t}\) be the split-site counts.  Rotating the
minus tail in both halves, including block phases and the twist, gives
\[
                    k^{1+(z_{l,+}-z_{r,+}+z_{r,-}-z_{l,-})/2}.
\]
The tail exponent is twice the number of relative plus roots in its
\(r\)-list minus twice the number of relative minus roots in its
\(l\)-list; the twist exponent is \(-2(m-1)\).  The split and cup
phases contribute \(k\) per group in \(\mathrm A,\mathrm B_0\),
and 1 in \(\mathrm B_1\), leaving the constant \(k^2\).

For the remaining pair comparison, record phases in units of \(\pi/8\)
modulo 16.  Temporarily omit the cup numerator and own-probe residues.
Remove the baseline ordinary cross-group sign
\((-1)^{1-(a-1)(b'-1)}\), including changed sites; its total is
sector-dependent only.  Matching spectator pairs now have ratio 1.
A full plus exceptional list of type \(\mathrm A\) or
\(\mathrm B_0\), against a spectator of length \(j\), has phases
\[
\begin{array}{c|ccc}
 j&0&1&2\\ \hline
 \text{same group}&12&8&12\\
 \text{opposite group}&4&8&4
\end{array}.
\]
Negate entries for \(\mathrm B_1\), and negate them for a minus
exceptional list, combining both negations when applicable.  The
plus-list versus minus-list phases in order
\(\mathrm A,\mathrm B_0,\mathrm B_1\), and the phases within one
exceptional list, are
\[
 \begin{pmatrix}0&4&0\\12&0&4\\0&12&0\end{pmatrix},
                         \qquad (0,4t,-4t).
\]
For example, place a plus exceptional list of type \(\mathrm A\)
before a same-group spectator of string length zero, and write \(x\)
for mark minus spectator.  The mark has labels \(l=1,r=0\), whereas
the spectator has \(l=r=-1\).  The \(l\)-half contributes 1.
The two mark roots in the \(r\)-half, including the single merge
divisor, give
\[
 \frac{P_{-3}(x)P_{-4}(x)}{P_{-1}(x)P_0(x)}
 \frac{P_{-1}(x)P_{-2}(x)}{2\cosh(x/2)}
 =i\frac{P_{-2}(x)P_{-3}(x)}{P_0(x)}.
\]
The corresponding ordinary-mark and probe factors on the gas side are
\[
 \frac{P_2(-x)P_3(-x)}{P_0(-x)}
 =-\frac{P_{-2}(x)P_{-3}(x)}{P_0(x)}.
\]
Their ratio is \(-i\), the phase 12 in the table.  Here we used
\(P_{-4}(x)=2i\cosh(x/2)\) and \(P_j(-x)=-P_{-j}(x)\).
The other entries use the same root and merge rules.

The following substitution verifies these tables and the constancy of
the ratios.  Encode powers of \(P_j\) by \(\xi^j\) in the signed
spin-minus-gas polynomial of each pair.  Insert the preceding root
factors at shift \(g-h\), subtract \(\xi^4\) at each merge shift,
and subtract \(\xi^{-ta}\) for the special denominator.  Subtract
ordinary string exponents for two ordinary sites.  For an ordinary
site of group \(z\), length \(a\), against a mark, subtract at
ordinary minus mark
\[
 \xi^{2z}+\xi^{3z}
                  +\sum_{h=1}^a\xi^{zh}(\xi^{-4z}-\xi^z).
\]
For a probe shifted from its site by \(b=0\) in \(\mathrm A\), or
\(b=ta\) in \(\mathrm B_a\), add for another ordinary site of
group \(z\), length \(j\), at ordinary minus probe-site,
\[
 \xi^{z-t-b}\left[1-
       \sum_{h=1}^j\xi^{zh}(\xi+\xi^{-1}-1)\right].
\]
For another mark add, at probe-site minus mark,
\(\xi^b(1+\xi^{2t})\) in the same group and \(-\xi^{b-3t}\) in
the other group.  Reverse exponents for reversed arguments.  Substituting
the displayed \(\mathcal R\)'s gives zero modulo \(\xi^8-1\).
To recover the constant send the ordered real difference to \(+\infty\):
sum coefficients times exponent plus 8, omit the added 8 on reversed
probe and ordinary-mark monomials, and replace it by \(-4\) on merge
monomials.  This gives the displayed tables after baseline removal.

The residues of the singular own-probe factors are
\(-1/P_{2t}(0),1,P_t(0)\) for
\(\mathrm A,\mathrm B_0,\mathrm B_1\); the remaining factors
are evaluated at the pole.  Write these residues as the vector $R$.
The exceptional scalar ratios relative to $C_0$ are $E=(1,1,2c)$,
and the cup numerators are $N=(1,-2itd,-2itd)$.  In the
spin-to-residue ratio their combined contribution is
$EN/R=(2itd,-2itd,2d)$.  They have equal modulus
\(2d\), and relative phases \((0,8,-4t)\), since
\(\sin\lambda=c\).  With the internal exceptional
phases this is \((0,8+4t,-8t)\).  Total ordinary length \(n-1\)
makes the spectator contribution zero if exactly one \(\mathrm B_1\)
occurs, and \(8(1+\#\mathrm B_0)\) otherwise.  Adding the matrix
above gives 8 in all nine cases.  Thus the ratio is independent of pole
sites and string lengths within a sector.  Finally, replacing ordinary
baselines by \(2d\) times their internal residues contributes only
\((c/d)^{\sum(1-a)}=c/d\).  All remaining magnitudes are fixed
constants, proving the identity and uniform bound.
\end{proof}

\subsection{Neutralizing the two probes}

Set $L_X=\gamma^{-1}\log s$ and $L_Y=\gamma^{-1}\log b$.
From now on the real site offsets are zero.  The following estimate
is the single-arc output needed in geometric applications.  The
statement includes the strongly unequal-wall regime.

\begin{theorem}[Uniform one-arc bound]
\label{09_fast_paths:one-arc-bound}\label{thm:one-arc}
There are constants $\delta_1>0$, $L_*>0$, and $C<\infty$ such
that
\[
                         A(p,p')\le b^C
\]
whenever $L_Y\ge L_X\ge L_*$, $s\le b$, and at least
$1/2-\delta_1$ of the ordinary $X$ sites have height $\beta$.
The marks may each have either height in $\{0,\beta\}$.
There is no positive lower bound on $L_X/L_Y$.
\end{theorem}

The proof has three tasks.  We move the probe integral to the positive
magnetic mean contours, keeping the residues that bind a particle to a
probe.  We then add two positive half-masses to make the fluctuation
field neutral.  Coercivity and the real reference comparison finally
bound the remaining integral by $\exp(O(L_Y))$.  The last step below
checks uniformity as $L_X/L_Y$ tends to zero.

\begin{proof}
Write $L=L_Y$.  Initially suppose that $L_X\ge c_1L$ with fixed
$c_1>0$ and let $L\to\infty$.  We will identify all uses of this
temporary restriction and remove them at the end.

\paragraph{Deformation and bound particles.}
Move the doubles to the final magnetic lines.  Their fused factors
cancel the probe denominators, so this movement crosses no new pole.
Move the minus far singleton from height $-2.35\beta$ to
$-7\beta/2$.  The only crossed poles bind a particle to the upper
edge of the plus probe, giving particle height $-2.66\beta$, or
to the lower edge of the minus probe, at height $-2.72\beta$.
A pair zero prevents two particles from binding to the same probe.
After a binding, the distances swept from that particle during the
remaining minus far-singleton movement are strictly less than one in
units of $\beta$.  Move free variables
of the same label together; this avoids any additional particle pole.
There are therefore only boundedly many binding types, each with at
most two attached particles.

The factorial measures survive this operation exactly.  If a species
has $r$ particles, summing the $r$ possible choices of the attached
particle changes $1/r!$ into $1/(r-1)!$.  For two distinct choices,
$r(r-1)/r!=1/(r-2)!$.  The same cancellation applies when the
particles come from different species.  Thus the unbound variables
retain one factorial divisor for each species, with no extra growing
factor from the labels of the bound particles.

Move the main $X$ line to the positive-density magnetic graph.  On
the final contours all denominators are regular after fusing doubles
and deleting each binding denominator.  The joint tail exponent is
bounded above by
\[
                  -\sum_v|\Re v|-\sum_{t=\pm1}|\Re y_t|
\]
up to a constant depending on the counts.  This proves simultaneous
absolute convergence and justifies the deformations and the limit of
coinciding real site offsets.

We also need a small extension below upper-site fraction $1/2$.
For any sufficiently small fixed $\delta_2>0$, decrease $\delta_1$
so that the positive-density graph has heights in
$[1/2-\delta_2,1]$ in units of $\beta$.  The positive-real-part
implicit-function argument of Lemma~\ref{lem:06-magnetic-curve}
continues uniformly on this slightly larger parameter interval.
All regular offsets retain strict margins.  The boundary-line
positivity used in Theorem~\ref{thm:06-real-positivity} persists by
compact continuity of its normalized symbol.  At infinity this
perturbation creates no new offset-zero coincidence, and the limiting
off-diagonal pair remains unchanged.  Consequently the magnetic mean
and real-energy estimates used below continue with uniform constants.

\paragraph{Mean cancellation and neutralization.}
Fix a binding type and freeze the probes.  The ordinary mean and
marked-anomaly identities of Lemma~\ref{lem:06-magnetic-mean}
apply to attached particles as well, because their shifts remain
admissible.  Each probe has an additional exact real mean cancellation:
its ordinary-site log is the negative of its particle log integrated
against $\mu$.

Here is the transform check, including the branch information.
Evaluate the mean on group $\sigma$ horizontally.  Main height minus
probe height plus $\sigma-t$ belongs strictly to $(1,7)$ or
$(-7,-1)$, and the corresponding site difference stays in the same
component of $\mathbb R\setminus8\mathbb Z$.  The three probe
shifts $+1,-1,0$, with coefficients $+1,+1,-1$, multiply the
logarithmic transform by
$c_*(k)=2\cosh(\beta k)-1$.  Each translated copy of the base density $\rho$ contributes its
translation factor times $1/c_*(k)$.  Multiplying by $c_*(k)$
cancels the density transform while retaining that site translation;
summing the upper and lower copies gives exactly the site transform,
with no angle fold.  After removing the common linear growth, the exponential
tails fix the real constant.  The remaining mark factor satisfies
\begin{equation}
                 \log|K_t|\le C-\tfrac12|\Re y_t|.
\label{eq:one-arc-probe-tail}
\end{equation}

Split main particles at their own walls $[-L_g,L_g]$.
Let $\alpha$ denote the interior trial densities and $\zeta$ the
positive counting measures of the remaining free particles together
with the attached particles.  In the reference real form, place an
attached particle on the standard minus far-singleton line; its actual
line will be accounted for by a localized error.  Let $\psi$ be the
even positive smoothing probability already used in
Lemma~\ref{lem:07_magnetic_analysis-smoothing}.  On that same minus
far-singleton label define
\[
 \chi_t(u)=\tfrac12\psi(u-\Re y_t),\qquad t=\pm1,
 \qquad
 x=\Re\alpha-\mu+\psi*\zeta+\chi_++\chi_-.
\]
All other components of $\chi_t$ are zero.  The one-arc gas has
$n-1$ ordinary elementary particles, whereas $\mu$ has total mass
$n$.  Since each $\chi_t$ is nonnegative and has mass $1/2$, the
weighted total mass of $x$ is zero.  Thus $B(x)$ is defined and the
magnetic coercivity theorem applies.  More precisely, put $M=1-m$ in the gas identity.  The $X$ total
is $s-M$, and the $Y$ total after adding the two half-masses is
$b+M$.  Thus the count tests have exactly the magnetic group totals,
with only the bounded extra positive masses in their local estimates.

Let $\mathcal D_{\rm arc}$ be the centered trial cost after subtracting
the logarithms of the probes and their binding residues.  The extracted
site factor remains
\[
       \mathcal P_N=(\gamma/2)^N\prod_{i<j}H_*(a_i-a_j)^2
\]
up to a bounded multiplier.  The change by one in the ordinary total
and the at most two binding residues contribute only bounded activity
and scalar factors.  We claim that, for fixed sufficiently small
smoothing scale $h$,
\begin{equation}
 \Re\mathcal D_{\rm arc}
 \ge cB(x)+\tfrac12\sum_{t=\pm1}|\Re y_t|
                  -C_h(L+n_e),
\label{eq:09-trial}
\end{equation}
where $n_e$ is the number of explicit variables, including the
bounded number of attachments.

To verify this inequality, first retain only the linear tails.  After
mean cancellation, the probe contribution apart from $K_t$ is
\[
 -\tfrac12\sum_{t=\pm1}
 \left(\sum_Iw_I(\Re\alpha-\mu+\zeta)_I,
                              |\,\cdot-\Re y_t|\right).
\]
This is exactly the cross term with two point half-masses for the tail
kernel $-w_Iw_J|u-v|$.  Replacing those point masses by $\chi_t$
changes the expression by exponentially localized tests, because
$\psi$ is even.  The compensator self and mutual terms are bounded
above uniformly in the two centers: the negative linear mutual tail
has a favorable sign, and the smoothed local remainders are bounded.
A removed binding denominator changes no tail, since its binding
distance is zero.  Equation~\eqref{eq:one-arc-probe-tail} supplies
the remaining positive term in the trial cost.

It remains to control the localized parts of the kernels.  The residual
probe logs have uniformly bounded localized $H^1$ norms on the main
curves.  On explicit labels and residues they have exponentially
localized upper bounds.  Negative Sobolev duality from magnetic
coercivity bounds their pairing against the interior-minus-mean field
by an arbitrarily small positive fraction of $B(x)$ plus a constant.
In applying it, write that field as
$x-\psi*\zeta-\chi_+-\chi_-$; this retains the compensator terms.
Averaging far-label logs against $\psi$ gives the required smooth
compensator tests.  The explicit parts are controlled by bounded
localized kernels and the cell estimate
\[
 \sum_j n_j^2\le C(1+|\log h|)(B(x)+1).
\]
The added positive masses have bounded total, so this estimate and
the remaining count estimates of
Lemma~\ref{lem:07_magnetic_analysis-smoothing} have only bounded
additive changes.

For an attached particle, compare its actual line with the standard far
line only after smoothing.  Against an interior variable the tails
agree and their difference is a regular localized test.  Against
another explicit charge compare the actual action directly with the
twice-smoothed standard kernel.  Regularity of the actual denominators,
equality of their tails, and boundedness of the smoothed local kernel
show that the negative error is at most
$C_h e^{-c|u-v|}$.  In particular no comparison at an unsmoothed
fictitious collision is used.  Since there are at most two attachments,
summing these explicit-pair errors costs at most $C_h n_e$.
The localized tests paired against $x$ cost at most
$\eta B(x)+C_{h,\eta}$ by Sobolev duality, for any fixed $\eta>0$;
pairing against $\psi*\zeta$ costs another $C_h n_e$, and the two
compensators cost a bounded amount.  Thus these attachment errors do
not require a coefficient tending to zero with $h$.
The marked anomalies satisfy the same bounds.
Equation~\eqref{eq:M-22} controls free explicit pairs, with
an omitted-diagonal cost at most $C_h n_e$.  The remaining mixed
errors are paired against $x-\psi*\zeta-\chi_+-\chi_-$.
Choose $h$ after the fixed far-contour parameter so that the smoothing
coefficient is small enough to absorb the quadratic errors into
$B(x)$.  These estimates prove Equation~\eqref{eq:09-trial}.

\paragraph{Real reference integration and factorial summation.}
We now integrate the interior fluctuations with the probes and explicit
configuration frozen.  The interior Hessian and its diagonal compensation
are the same as in the magnetic gas without the probes.  The new linear
tests are smooth in the reflected reference coordinates at each fixed
size.  The magnetic decomposition also contains the attenuation
$\mathfrak a=\exp(-(X_{\rm int},R\zeta))\in[0,1]$, where
$X_{\rm int}$ is the interior empirical configuration and $R$ is
the nonnegative horizontal cutoff kernel.  Take the absolute value
of the integral before replacing this factor by one; no signed or
oscillatory identity is inferred by discarding it.  Apply the coarse
real source estimate
\eqref{eq:G-13} and the compensated quadratic integration estimate
\eqref{eq:G-18}, with zero displacement.  The latter evaluates only
the smooth compensated cost on an empirical field, not its two singular
energies separately.

For clarity, the spectral loss is $O(L)$ already at this coarse level.
Let $F$ be the real Hessian relative to $h_0$.  The magnetic verification
of Equation~\eqref{eq:G-8} and the trace bound give
\[
 F\ge cL^{-2}I,\qquad \|F-I\|_1\le CL,
 \qquad N_F(a)\le CL\sqrt a\quad(0<a<a_0).
\]
On $[a_0,1)$, $-\log a\le C_{a_0}(1-a)$, so the trace norm
bounds that part of the negative logarithmic sum by $CL$.  Below
$a_0$, the sum is
\[
 (-\log a_0)N_F(a_0)+\int_0^{a_0}\frac{N_F(a)}a\,da,
\]
which is again $O(L)$ because
$\int_0^{a_0}a^{-1/2}\,da<\infty$.  Thus
\[
                  \sum_{\lambda<1}-\log\lambda\le CL.
\]
The unnormalized reference masses must also be retained.  For each
nonnegative integer dimension $N_g$ they satisfy
$J_{N_g}=\exp(-\tfrac14\log(N_g+1)+O(1))$, with
$J_0=1$, and hence are bounded above.  All dimensions are fixed
integers during this real integration and at most $n-1\le 2b$.
Consequently the reference integral is at most the exponential of
minus its real trial infimum plus $CL$.  This use of the source
estimate needs neither a count truncation nor the sharp magnetic
transport theorem.

Reserve a fixed fraction of $B(x)$ in
Equation~\eqref{eq:09-trial}.  The magnetic count bounds, with the
bounded extra masses, give
\[
 \ell_e+\frac{M^2+n_e^2}{L}\le C_h(B(x)+1),
 \qquad
 \ell_e=\sum_{v\in\zeta}(|\Re v|-L)_+.
\]
The inequality
$C_hn_e\le \eta n_e^2/L+C_{h,\eta}L$, with suitably small
$\eta>0$, absorbs the linear cost in a further reserved fraction.
We obtain positive damping in $\ell_e$, $M^2/L$, and $n_e^2/L$,
as well as $\tfrac12\sum_t|\Re y_t|$, at a cost $CL$.

Here is the resulting absolute summation.  Each free explicit position
has an exponentially damped tail outside $[-L,L]$, so its integral,
including its fixed activity, is at most $C(L+1)$.  If the free counts
are $r_I$, their factorial sum is bounded by
\[
 \sum_{(r_I)\ge0}\prod_I\frac{[C(L+1)]^{r_I}}{r_I!}
                     =\exp(O(L)).
\]
This enlargement drops the favorable quadratic count damping and the
finite-degree restriction.  Summing the remaining imbalance damping
$\exp(-cM^2/L)$ over $M\in\mathbb Z$ costs $O(\sqrt L)$.
The at most two attachments have already been absorbed into the probe
integrals and their factorial cancellation.  Finally
\[
 \prod_{t=\pm1}\int_{\mathbb R}e^{-|v_t|/2}\,dv_t<\infty.
\]
The bounded number of contour edges and binding types adds a fixed
factor.  Therefore the absolute contour sum divided by
$|\mathcal P_N|$ is at most $\exp(CL)$.

The uniform denominator theorem, Theorem~\ref{thm:06-denominator},
applies to all $N$ site heights in $\{0,\beta\}$ without a
restriction on their two multiplicities.  Dividing the last estimate
by $|\widetilde h_N^2/\mathcal P_N|$ proves $A(p,p')\le b^C$
in the temporary comparable-wall regime.

\paragraph{Uniformity for arbitrarily unequal walls.}
We used only the count and smoothing estimates
\eqref{eq:M-20}--\eqref{eq:M-23}, an $O(L_Y)$ trace bound,
and the gap and small-value count in Equation~\eqref{eq:G-8}.
We did not use the sharp determinant asymptotic or perform a physical
shift.  These coarse ingredients extend as follows.

Keep a separate unstretched reference coordinate on each wall.  Counts
and tails always use $L_Y$, while a grid for main explicit particles
starts at that label's own wall.  Denote the upper ordinary-site
fraction by $t_X$.  If $t_X$ tends to one without being identically
one, put
$d_0=\beta(1-\theta(0))$.  Its discreteness gives
$1-t_X\ge1/s$ whenever it is nonzero.  The curve estimates therefore
give $1+|\log d_0|=O(L_X)$.  The exceptional localized
half-derivative loss at the curved origin is only polynomial in this
quantity.  A main explicit particle from that label lies beyond its
own wall; smoothing it into the exceptional neighborhood gains an
exponential buffer of order $L_X$.  Thus the smoothing errors formerly
written $o_L(1)$ become $o_{L_X}(1)$, with no wall-ratio assumption.
All other offset, norm, and cell estimates are uniform on their separate
wall intervals.

The regular trace comparison has bounded cost per unit cell, hence
total cost $O(L_Y)$.  At the rare center, the dyadic comparison in
Equation~\eqref{eq:M-32} uses
$O(1+\log(1/d_0))=O(L_X)$ shells.  Its coarse trace cost is
therefore still $O(L_Y)$.  The caps and the block-mean proof of the
gap and small-value count are unchanged: their two interval lengths
are bounded above by $L_Y$.  This proves precisely the spectral
estimates used in the real comparison.  Each fixed-size kernel and
linear test is smooth, which is all that the coarse source estimate
requires.  No lower bound for the rare central density and no polynomial
margin for its derivatives is needed, because particles are never
shifted.

Fix the far parameter and then a sufficiently small smoothing scale.
The bounds just proved give a bounded exponent along every sequence
with $L_Y\ge L_X\to\infty$, taking convergent mixture
subsequences when needed.  If no fixed $L_*$ and uniform exponent
existed, one could choose examples with $L_X\ge j$ and
$\log A(p,p')/\log b>j$.  A subsequence would contradict these
same coarse bounds.  Hence some fixed $L_*$ and $C$ work throughout
the stated domain, completing the proof.
\end{proof}
